\begin{thebibliography}{99}
\addcontentsline{toc}{chapter}{参考文献}
\small

\bibitem{boyd2004}
S. Boyd and L. Vandenberghe.
\emph{Convex Optimization}. Cambridge University Press, 2004.
\href{https://web.stanford.edu/~boyd/cvxbook/}{作者提供的书籍与配套资源}。

\bibitem{cvxexamples}
CVX Research and example contributors.
\emph{CVX Example Library}, 包含 A. Mutapcic 编写的功率控制与阵列设计示例，2006。
本书核查版本：\href{https://github.com/cvxr/CVX/tree/a0a775908bc1eec1625d56498d0224132b0d8b7f/examples}{CVX commit a0a7759}。
\href{https://github.com/cvxr/CVX/blob/a0a775908bc1eec1625d56498d0224132b0d8b7f/LICENSE.txt}{许可证中 Example library 小节}将示例库列为公有领域；核查日期 2026-11-08。

\bibitem{cvxpyexamples}
R. Gowers, R. Hill, S. Al-Izzi, T. Pollington, and K. Briggs.
\emph{Water Filling in Communications};
\emph{Power Assignment in a Wireless Communication System}.
CVXPY 官方示例，\href{https://github.com/cvxpy/examples/tree/73a6892ef5bea4ded68f8f9fab17e48119dc62f5/notebooks/WWW}{核查版本 commit 73a6892，2026}。
\href{https://github.com/cvxpy/examples/blob/73a6892ef5bea4ded68f8f9fab17e48119dc62f5/LICENSE}{Apache-2.0}；核查日期 2026-11-08。

\bibitem{tse2005}
D. Tse and P. Viswanath.
\emph{Fundamentals of Wireless Communication}.
Cambridge University Press, 2005.
\href{https://web.stanford.edu/~dntse/wireless_book.html}{作者提供的在线章节}。

\bibitem{yates1995}
R. D. Yates.
A framework for uplink power control in cellular radio systems.
\emph{IEEE Journal on Selected Areas in Communications}, 13(7):1341--1347, 1995.
\href{https://doi.org/10.1109/49.414651}{doi:10.1109/49.414651}.

\bibitem{kelly1998}
F. P. Kelly, A. K. Maulloo, and D. K. H. Tan.
Rate control for communication networks: shadow prices, proportional fairness and stability.
\emph{Journal of the Operational Research Society}, 49(3):237--252, 1998.
\href{https://www.statslab.cam.ac.uk/~fpk1/rate.html}{作者论文主页}。

\bibitem{lebret1997}
H. Lebret and S. Boyd.
Antenna array pattern synthesis via convex optimization.
\emph{IEEE Transactions on Signal Processing}, 45(3):526--532, 1997.
\href{https://web.stanford.edu/~boyd/papers/ant-pat.html}{作者论文主页}。

\bibitem{luo2010}
Z.-Q. Luo, W.-K. Ma, A. M.-C. So, Y. Ye, and S. Zhang.
Semidefinite relaxation of quadratic optimization problems.
\emph{IEEE Signal Processing Magazine}, 27(3):20--34, 2010.
\href{https://doi.org/10.1109/MSP.2010.936019}{doi:10.1109/MSP.2010.936019}.

\bibitem{ben2009}
A. Ben-Tal, L. El Ghaoui, and A. Nemirovski.
\emph{Robust Optimization}.
Princeton University Press, 2009.
\href{https://press.princeton.edu/books/ebook/9781400831050/robust-optimization-pdf}{出版社书籍页面}。

\bibitem{shi2011}
Q. Shi, M. Razaviyayn, Z.-Q. Luo, and C. He.
An iteratively weighted MMSE approach to distributed sum-utility maximization for a MIMO interfering broadcast channel.
\emph{IEEE Transactions on Signal Processing}, 59(9):4331--4340, 2011.
\href{https://doi.org/10.1109/TSP.2011.2147784}{doi:10.1109/TSP.2011.2147784}.

\bibitem{shen2018}
K. Shen and W. Yu.
Fractional programming for communication systems---Part I: Power control and beamforming.
\emph{IEEE Transactions on Signal Processing}, 66(10):2616--2630, 2018.
\href{https://doi.org/10.1109/TSP.2018.2812733}{doi:10.1109/TSP.2018.2812733}.

\bibitem{dinkelbach1967}
W. Dinkelbach.
On nonlinear fractional programming.
\emph{Management Science}, 13(7):492--498, 1967.
\href{https://doi.org/10.1287/mnsc.13.7.492}{doi:10.1287/mnsc.13.7.492}.

\bibitem{razaviyayn2013}
M. Razaviyayn, M. Hong, and Z.-Q. Luo.
A unified convergence analysis of block successive minimization methods for nonsmooth optimization.
\emph{SIAM Journal on Optimization}, 23(2):1126--1153, 2013.
\href{https://doi.org/10.1137/120891009}{doi:10.1137/120891009}.

\bibitem{beck2009}
A. Beck and M. Teboulle.
A fast iterative shrinkage-thresholding algorithm for linear inverse problems.
\emph{SIAM Journal on Imaging Sciences}, 2(1):183--202, 2009.
\href{https://doi.org/10.1137/080716542}{doi:10.1137/080716542}.

\bibitem{boyd2011}
S. Boyd, N. Parikh, E. Chu, B. Peleato, and J. Eckstein.
Distributed optimization and statistical learning via the alternating direction method of multipliers.
\emph{Foundations and Trends in Machine Learning}, 3(1):1--122, 2011.
\href{https://web.stanford.edu/~boyd/papers/admm_distr_stats.html}{作者论文与配套实现主页}。

\bibitem{neely2010}
M. J. Neely.
\emph{Stochastic Network Optimization with Application to Communication and Queueing Systems}.
Morgan \& Claypool, 2010.
\href{https://doi.org/10.2200/S00271ED1V01Y201006CNT007}{doi:10.2200/S00271ED1V01Y201006CNT007}.

\bibitem{writing2026}
hongjinrui.
\emph{科研写作（research-writing）规则}, 2026.
\href{https://github.com/hongjinrui/writing_skills/blob/main/SKILL.md}{SKILL.md}。
本书采用其符号定义、推导衔接、假设与证据、图表可复现等写作规则；读取日期 2026-11-08。

\bibitem{shapiro2009}
A. Shapiro, D. Dentcheva, and A. Ruszczy\'nski.
\emph{Lectures on Stochastic Programming: Modeling and Theory}.
SIAM and MPS, 2009.
\href{https://doi.org/10.1137/1.9780898718751}{doi:10.1137/1.9780898718751}.

\bibitem{puterman1994}
M. L. Puterman.
\emph{Markov Decision Processes: Discrete Stochastic Dynamic Programming}.
Wiley, 1994.
\href{https://doi.org/10.1002/9780470316887}{doi:10.1002/9780470316887}.
\end{thebibliography}
